\chapter{How to Use This Book}

No previous knowledge is assumed. Readers who have never programmed, or have
never looked inside a computer, should begin with the chapter \emph{Starting
from Zero}, on page~\pageref{ch:zero}: a few pages that explain the words and
the ideas used from the first lecture on.

\section*{One chapter, one lecture}

The book has twenty-five chapters, numbered from~0 to~24 as the lectures of the
course are, and as almost everything else in this book is: bits, registers,
rows and columns all start from zero. Chapter~$N$ follows Lecture~$N$, which
has two slide decks. The \emph{lecture deck}, of at most twenty slides, holds
the concepts; the \emph{code deck} holds everything that is typed and run:
each listing with its explanation, the complete program of the lecture, what
it shows, and the way to test it on \prgqt, on QEMU and on the board. The
chapter is the long form of both: the order of the sections is the order of
the slides, the listings are the same files, the questions at the end are
those shown in class, and what twenty slides have no room for is here. A
reader who has attended the lecture finds the sentences that a slide cannot
hold; a reader who has missed it finds the lecture.

Most chapters end with a section, \emph{The program of this chapter}, that
prints in full the program from which the listings of the chapter are taken,
shows the frame that it produces and gives the commands to run it. The
programs are in the folder \fname{code} of the repository of the course, one
per lecture, and every one of them is built and run by the automatic tests;
Section~\ref{sec:prg:testing} is the step-by-step guide for the three hosts.

\begin{center}\small
\begin{tabularx}{\textwidth}{@{}lP{0.36\textwidth}L@{}}
\toprule
\textbf{Part} & \textbf{Chapters} & \textbf{Then, the project}\\
\midrule
Start & 0 Introduction, 1 Setting up, 2 First contact & \\
Information & 3 History, 4 Digital information, 5 Integers, 6 Reals & 7 \vermone, part 1\\
Logic & 8--9 Boolean algebra, 10 Electronics, 11 Combinational machines & 12 \vermone, part 2\\
State and memory & 13 Sequential machines, 14 Registers and buses, 15 Memory, 16 Permanent memory & 17 \vermone, part 3\\
The machine & 18 Processor, 19 Instructions and interrupts, 20 Video card, 21 Input and output & 22 \vermone, part 4\\
Publication, outlook & 24 Beyond one core & 23 \vermone on the Cartridge Store\\
\bottomrule
\end{tabularx}
\end{center}

\section*{The shape of a topical chapter}

Every chapter that presents a topic has the same five movements.

\begin{enumerate}
\item \textbf{Warm-up on the machine.} A few lines of \riscv on \prg that
  produce a fact needing an explanation. Run them before reading on.
\item \textbf{The formal treatment.} Definitions, theorems, methods and worked
  examples.
\item \textbf{Back to the machine.} The same rules, as \riscv implements them.
\item \textbf{Check your understanding.} Multiple-choice questions in the style
  of the midterm test. Answers, with a line of explanation each, are in
  Appendix~\ref{app:answers}: read them after answering, not instead.
\item \textbf{Exercises, glossary, reading.} Exercises are done on paper first
  and on \prg afterwards; the glossary gives the Italian term for every English
  one, since the course is taught in a bilingual setting.
\end{enumerate}

\section*{The shape of a tutorial chapter}

Chapters 7, 12, 17, 22 and 23 build \vermone. They are self-contained and
meant to be followed at the keyboard. Each step shows a few lines, says
whether they are \textbf{new}, \textbf{changed} or \textbf{unchanged}, and ends
with a \emph{Run it} checkpoint stating what the screen must show. If a
program refuses to work at the beginning of a tutorial, start from the
reference file of the previous part (Appendix~\ref{app:verm1}): nobody should
be left behind by a bug.

\section*{Boxes}

\begin{prgconcept}
A \textbf{Concept} box states an idea that the rest of the chapter depends on.
\end{prgconcept}
\begin{definition}{how definitions look}
Formal definitions and theorems are numbered within the chapter.
\end{definition}
\begin{prgtry}
A \textbf{Try It Yourself} box asks for something to be done at the keyboard,
and says what must be observed.
\end{prgtry}
\begin{prgwarn}
A \textbf{Watch Out} box marks a mistake that is made every year.
\end{prgwarn}
\begin{prgmachine}
A \textbf{Back to the Machine} box ties a formal result to an instruction, a
register or a screen of the debugger.
\end{prgmachine}
\begin{prgaside}
A \textbf{Historical Note} box tells where an idea comes from.
\end{prgaside}

\section*{Learning status: understand now, use now, understand later}

The course programs first and formalises later, so a lecture always shows more than it
explains. Each chapter therefore opens, after its goals, with a box that sorts what it
contains into three states, and the slides carry the same badges above their titles.

\begin{description}[style=nextline,leftmargin=1.5em]
\item[\masterychip{now}{UNDERSTAND NOW}] By the end of the lecture you can explain it,
  trace it by hand and apply it on your own. The questions and the exercises check these.
\item[\masterychip{use}{USE NOW}] You follow and use the recipe today: a directive, a command,
  a call of the runtime, a register convention. How it works inside is explained in the
  chapter given in brackets, or is not a subject of the course (the tools of \prg are used,
  not studied).
\item[\masterychip{later}{UNDERSTAND LATER}] A preview: recognise the name and what it is
  for. The chapter given in brackets explains it; nothing is deferred without a destination.
\end{description}

The states are expectations in time, not levels of difficulty: the same idea often moves
from \emph{later} to \emph{use} to \emph{now} over a few chapters (the calling convention is
used from Chapter~\ref{ch:firstcontact} and studied in Chapter~\ref{ch:processor}). A slide or a
section without a badge is context: a story, an example, a question, an exercise. Colour is
never the only cue: the book, the wrench and the clock, and the solid, double and dashed
frames, tell the three states apart in grey print too.

\section*{Notation}

Numbers without a mark are decimal. Hexadecimal numbers are written as the
assembler wants them, \hex{07E0}; binary numbers as \code{0b1010} in code and
with a subscript, $1010_2$, in formulas. Bits are numbered from the right
starting at zero, so that bit~$i$ has weight $2^{i}$. Registers
(\reg{a0}, \reg{sp}), instructions (\ins{addi}) and file names are set in a
typewriter face. In Boolean expressions the complement of $x$ is $\nott{x}$,
the sum is $+$ and the product is written by juxtaposition. Multiples of bytes
follow the binary prefixes where precision matters:
$1~\mathrm{KiB}=2^{10}$ bytes, $1~\mathrm{MiB}=2^{20}$ bytes.

\section*{The appendices}

Appendix~\ref{app:riscv} is a complete reference of the \riscv instruction set
used in the book: registers, formats, every instruction of RV32I and of the M,
A, C and Zicsr extensions, pseudo-instructions, assembler directives, the
calling convention and the machine-mode control registers.
Appendix~\ref{app:prg32} is the reference of \prg: the contract, the
programming interface, the cartridge format, the command-line tools, the hosts
and the Cartridge Store. Appendix~\ref{app:github} is a reference of Git and
GitHub, the tools with which the platform, this book and the reader's project
are kept and shared. Appendix~\ref{app:verm1} prints the four reference
programs of \vermone in full; Appendix~\ref{app:answers} the answers to all the
questions; Appendix~\ref{app:glossary} the complete English--Italian glossary.
